\begin{textsample}{POS Dim 5 – global\_north\_2023\_11 – Score 9.00 – global\_north\_2023\_11/103250287.txt}  \label{ex:f5_pos_416}
SERIES \&amp ; MORE POLITICO Live Follow us U.S. diplomats slam Israel policy in leaked memo The document offers a window into internal fury at the State Department over President Joe Biden ' s Middle East policies.
The \textbf{arguments} in the dissent memo offer a window into the thinking of many people at the State Department, which has long been vexed by the decades-old Israeli-Palestinian conflict.
Alastair Pike/AFP/Getty Images State Department staffers offered a blistering critique of the Biden administration ' s handling of the Israel-Hamas war in a dissent memo obtained by POLITICO, arguing that, among other things, the U.S. should be willing to publicly criticize the Israelis.
The message suggests a growing loss of confidence among U.S. diplomats in President Joe Biden ' s approach to the Middle East crisis.
It reflects the sentiments of many U.S. diplomats, especially at mid-level and lower ranks, according to conversations with several department staffers as well as other reports.
If such internal disagreements intensify, it could make it harder for the Biden administration to craft policy toward the region.
The memo has @ @ @ @ @ @ @ @ @ @, and that it balance its private and public messaging toward Israel, including airing criticisms of Israeli military tactics and treatment of Palestinians that the U.S. generally prefers to keep private.
The gap between America ' s private and public messaging"contributes to regional public perceptions that the United States is a biased and dishonest actor, which at best does not advance, and at worst harms, U.S. interests worldwide,"the document states."We must publicly criticize Israel ' s \textbf{violations} of \textbf{international} norms such as failure to limit offensive operations to legitimate military targets,"the message also states."When Israel supports settler violence and illegal land seizures or employs excessive use of force against Palestinians, we must communicate publicly that this goes against our American values so that Israel does not \textbf{act} with impunity."The memo is marked"sensitive but unclassified."It ' s not clear how many people signed it or if and when it was \textbf{submitted} to the department ' s Dissent Channel, where employees can voice policy @ @ @ @ @ @ @ @ @ @ was revised in any way beyond the version obtained by POLITICO.
Still, the \textbf{arguments} in it offer a window into the thinking of many people at the State Department, which has long been vexed by the decades-old Israeli-Palestinian conflict.
The department declined to comment directly on the memo, as is standard on such communications.
It referred POLITICO to past statements by spokesperson Matthew Miller, who has said Secretary of State Antony Blinken welcomes such \textbf{arguments} and weighs them carefully."One of the strengths of this department is that we do have people with different opinions,"Miller said about such messages during a press briefing last month."We encourage them to make their opinions known."Multiple dissent memos about this war are being circulated in the State Department in efforts to gather signatures.
These communications may or may not be \textbf{classified}, but their contents are rarely leaked.
The department ' s Dissent Channel is a long-established vehicle that allows staffers to freely express their discontent on a policy matter without @ @ @ @ @ @ @ @ @ @ was authored by two midlevel staffers who have worked in the Middle East, said a department employee who has seen the document and was granted anonymity to discuss a sensitive topic.
The memo concedes that Israel has a"legitimate right and \textbf{obligation}"to seek justice against the Palestinian militants of Hamas, who killed some 1,400 Israelis in a shocking Oct. 7 attack.
But it argues that"the extent of human lives lost thus far is unacceptable"-- referring to the thousands of Palestinians, most of them civilians and many children, killed by Israel in the days since.
The U.S."tolerance"for such a high civilian death toll"engenders doubt in the rules-based \textbf{international} order that we have long championed,"the document states.
It argues that the United States must hold both Israel and Hamas responsible for their actions.
The memo ' s demands are unlikely to get far with Biden or his top aides, at least not anytime soon.
The president, Blinken and others have ruled @ @ @ @ @ @ @ @ @ @ dismantle Hamas, which is based in the Gaza Strip.
Hamas embeds its arsenal and fighters throughout the civilian population, making it hard for Israel to delineate targets.
Israeli officials say they try to minimize civilian deaths but that a certain number are inevitable given how Hamas positions its people and assets.
The Biden team has increasingly shifted its public messaging to emphasize the importance of \textbf{safeguarding} civilians and following \textbf{international} \textbf{law}.
But it has largely avoided direct public criticism of Israeli actions.
Blinken has been holding listening sessions with groups of staffers unhappy with the trajectory of U.S. policy in the last few weeks.

% matched lemmas: act, argument, classify, international, law, obligation, safeguard, submit, violation
\end{textsample}
